\scanpage{15}
\deleted{d'ordre fini dans le groupe envisagé \inserted{\deleted{\unsure{stable}} \unsure{fixé} par $\unsure{\rho|\varphi}$}, il faudrait donc avoir
\[
\varepsilon_0\varepsilon_1\varepsilon_\infty=1
\]
i.e. en termes des $l_i,\sigma_i$
\[
l_1\underbrace{\sigma_0l_0}_{\substack{\inserted{\unsure{l_1l_\infty|l_\infty}}\\\deleted{l_1^{-1}l_\infty^{-1}\sigma_0}}}l_\infty\underbrace{\sigma_1l_1}_{l_\infty^{-1}l_0^{-1}}l_0\underbrace{\sigma_\infty l_\infty}_{l_0^{-1}l_1^{-1}}=1
\]
\margin{i.e. $l_\infty^{-1}l_0^{-1}l_1^{-1}=1$\quad OK}
On a donc
\[
\boxed{\varepsilon_0\varepsilon_1\varepsilon_\infty=1}
\]
on peut en tirer de plus jolies expressions pour les $\sigma_i$
\[
\begin{aligned}
\sigma_0&=\varepsilon_\infty^{-1}\varepsilon_1\\
\sigma_1&=\varepsilon_0^{-1}\varepsilon_\infty\\
\sigma_\infty&=\varepsilon_1^{-1}\varepsilon_0
\end{aligned}
\]
donc aussi la relation inattendue
\[
\sigma_\infty\sigma_1\sigma_0=1
\]}
Sans doute
\[
\varepsilon_0\varepsilon_1\varepsilon_\infty=\sigma_1
\]
i.e.
\[
\begin{aligned}
\varepsilon_0&=\varepsilon_1\varepsilon_\infty\varepsilon_1^{-1}\\
\varepsilon_1&=\varepsilon_\infty\varepsilon_0\varepsilon_\infty^{-1}\\
\varepsilon_\infty&=\varepsilon_0\varepsilon_1\varepsilon_0^{-1}
\end{aligned}
\]
\figure{\begin{tikzpicture}\draw (0,0)--(.15,0);\draw (.25,0)--(.8,0);\end{tikzpicture}}

Critique des notations et de l'approche. Les éléments $l_i,\sigma_i,\varepsilon_i,\rho\ldots$ sont dans les groupes discrets, qu'il faudrait distinguer des groupes profinis --- de plus il faudrait pour les groupes \unsure{discrets} \deleted{utiliser} (et $\unsure{\bar m|=}$ pour les
\mspage{14}
